\section{Formeln zur $\Gamma$-Funktion\label{gamma:section:formulas}}
\kopfrechts{Formeln zur $\Gamma$-Funktion}

Für die $\Gamma$-Funktion existieren verschiedene Darstellungen und Näherungsverfahren.
Im Folgenden werden ausgewählte Formeln sowie ihre jeweiligen Anwendungsgebiete vorgestellt, wobei auf ausführliche mathematische Herleitungen verzichtet wird.
Abschliessend wird die Bedeutung der $\Gamma$-Funktion für die Berechnung des Kugelnvolumens aufgezeigt.

\subsection{Produktformel von Euler}

Auf der Suche nach der $\Gamma$-Funktion im 18. Jahrhundert hat Leonhard Euler unendliche Produkte verwendet.
Mit Hilfe des Bohr-Mollerup-Theorems (Abschnitt~\ref{gamma:section:BMT}) kann nachgewiesen werden, dass es sich bei
\index{Produktformel, Gamma-funktion}%
\index{Gamma-Funktion!Produktformel}%
\index{Gamma-Funktion!Euler-Formel}%
\begin{equation}
\label{gamma-eq:product-formula-euler}
    \Gamma_{\text{Euler}}(z) = \lim_{n \to \infty} \frac{n!\,n^z}{z(z+1)\cdots(z+n)}
\end{equation}
um die $\Gamma$-Funktion handelt.

Die erste Eigenschaft des Bohr-Mollerup-Theorem lässt sich durch Einsetzen von $z=1$ überprüfen:
\begin{equation}
\label{gamma-eq:euler-first-property}
    \begin{split}
    \Gamma_{\text{Euler}}(1)
        &= \lim_{n \to \infty} \frac{n!\,n}{1 \cdot 2 \cdots (n+1)} \\
        &= \lim_{n \to \infty} \frac{n!\,n}{(n+1)!} \\
        &= \lim_{n \to \infty} \frac{n}{n+1} \\
        &= 1.
    \end{split}
\end{equation}
Für die zweite Eigenschaft wird $z+1$ anstelle von $z$ eingesetzt. Durch Umformen erhält man
\begin{equation}
\label{gamma-eq:euler-recursive-property}
    \begin{split}
    \Gamma_{\text{Euler}}(z+1)
        &= \lim_{n \to \infty}
            \frac{n!\,n^{z+1}}{(z+1)(z+2)\cdots(z+n+1)} \\
        &= \lim_{n \to \infty}
            \biggl(
            \frac{n}{z+n+1}
            \cdot
            \frac{n!\,n^z}{(z+1)(z+2)\cdots(z+n)}
            \biggr) \\
        &= \lim_{n \to \infty}
            \frac{n}{z+n+1}
            \cdot z
            \frac{n!n^{z}}{z(z+1)(z+2)\cdots(z+n)} \\
        &= \lim_{n \to \infty} \frac{n}{z+n+1}
            \cdot z\Gamma_{\text{Euler}}(z) \\
        &= z\Gamma_{\text{Euler}}(z).
    \end{split}
\end{equation}
Damit erfüllt die Produktformel die Eigenschaften~\ref{gamma-eq:bm-1} und~\ref{gamma-eq:bm-2}.
Jetzt muss man noch zeigen, dass $\Gamma_{\text{Euler}}$ gemäss Eigenschaft~\ref{gamma-eq:bm-3} logarithmisch konvex ist.
Die Log-Konvexität der Produktdarstellung wird im Standardbeweis des Satzes von Bohr-Mollerup bewiesen \cite[Abschnitt 4.1.4]{buch:MatSeSspezielleFunktionen}.
Nach dem Satz von Bohr-Mollerup gilt $\Gamma_{\text{Euler}}(z) = \Gamma(z)$.

\subsection{Integraldarstellung}

Die Integraldarstellung
\index{Integraldarstellung, Gamma-Funktion}%
\index{Gamma-Funktion!Integraldarstellung}%
\begin{equation}
\label{gamma-eq:integral}
    \Gamma_{\text{Integral}}(z) = \int_0^\infty t^{z-1} e^{-t}\,dt,
    \qquad \operatorname{Re}(z)>0
\end{equation}
eignet sich besonders gut zur Untersuchung von Ableitungen. 
Auch $\Gamma_{\text{Intergal}}$ kann mit dem Eigenschaften vom Bohr-Mollerup-Theorem (\ref{gamma:section:BMT}) überprüft werden. 

Die erste Eigenschaft des Bohr-Mollerup-Theorems lässt sich durch Einsetzen von $z=1$ überprüfen:
\begin{equation}
\label{gamma-eq:integral-first-property}
    \begin{split}
    \Gamma_{\text{Integral}}(1)
        &= \int_0^\infty t^{1-1}e^{-t}\,dt \\
        &= \int_0^\infty e^{-t}\,dt \\
        &= \left[-e^{-t}\right]_0^\infty \\
        &= 1.
    \end{split}
\end{equation}
Für die zweite Eigenschaft wird $z+1$ anstelle von $z$ eingesetzt. Durch partielle Integration erhält man für $z>0$
\begin{equation}
\label{gamma-eq:integral-recursive-property}
    \begin{split}
    \Gamma_{\text{Integral}}(z+1)
        &= \int_0^\infty t^z e^{-t}\,dt \\
        &= \left[-t^z e^{-t}\right]_0^\infty
            + z\int_0^\infty t^{z-1}e^{-t}\,dt \\
        &= z\Gamma_{\text{Integral}}(z).
    \end{split}
\end{equation}
Damit erfüllt die Integraldarstellung die Eigenschaften~\ref{gamma-eq:bm-1} und~\ref{gamma-eq:bm-2}.

Mit der Ableitungsformel 
\begin{equation}
    \Gamma_{\text{Integral}}'(z)=\int_0^\infty t^{z-1}e^{-t}\log(t)\,dt.
\end{equation}
kann man auch die logarithmische Konvexität (\ref{gamma-eq:bm-3}) von $\Gamma_{\text{Integral}}(z)$ beweisen, wie in \cite[Abschnitt 4.1.5]{buch:MatSeSspezielleFunktionen} gezeigt.
Aus dem Satz von Bohr-Mollerup folgt dann, dass auch $\Gamma_{\text{Integral}}(z)=\Gamma(z)$.

\subsection{Lanczos-Approximation}

Die Lanczos-Approximation
\index{Gamma-Funktion!Lanczos-Approximation}%
\index{Lanczos-Approximation}%
\begin{equation}
\label{gamma-eq:lanczos}
    \Gamma_{\text{Lanczos}}(z) \approx \sqrt{2\pi}
    \biggl(z + g - \frac{1}{2}\biggr)^{z - \frac{1}{2}}
    e^{-(z + g - \frac{1}{2})}
    \biggl(c_0 + \sum_{k=1}^{N} \frac{c_k}{z + k - 1}\biggr)
\end{equation}
wurde 1964 von Cornelius Lanczos vorgestellt und wird heute häufig in numerischen Bibliotheken eingesetzt.
\index{Lanczos, Cornelius}%
Der Parameter $g$ und die Koeffizienten $c_k$ werden dabei im Voraus festgelegt und müssen nicht bei jeder Auswertung neu berechnet werden.
Da nur eine endliche Summe mit $N$ Gliedern ausgewertet wird, ermöglicht das Verfahren eine genaue und zugleich effiziente Berechnung der $\Gamma$-Funktion. 
Auch hier gilt $\Gamma_{\text{Lanczos}}(z) = \Gamma(z)$. 
Eine ausführliche Herleitung der Approximation sowie eine Analyse ihrer Genauigkeit finden sich in \cite{LanczosGammaApproximation}.
